\chapter*{Appendix E: The Five Steps of Risk Assessment — Reference Card}
\label{chap:Appendix E: The Five Steps of Risk Assessment — Reference Card}

\paragraph{A concise reference guide to the five-step risk assessment process as described in Part III of this book.}

%---------------------------- SECTION ----------------------------
\section*{Step 1 — Identify Hazards and Risk Sources}
\paragraph{\textit{What are we looking for?} All sources of potential harm — conditions, activities, events, or circumstances that could cause adverse consequences if not adequately controlled.}
\paragraph{\textit{Key techniques:}}
\begin{itemize}
  \item{Process walkthrough and observation.}
  \item{Structured brainstorming (SWIFT, HAZOP).}
  \item{Historical incident and near-miss analysis.}
  \item{Bow-tie analysis.}
  \item{Regulatory and external hazard databases.}
  \item{Horizon scanning and environmental scanning.}
\end{itemize}

\paragraph{\textit{Key questions:}}
\begin{itemize}
  \item{What could go wrong?}
  \item{What conditions or activities create the risk?}
  \item{What external events could affect us?}
  \item{What have we missed?}
\end{itemize}

%---------------------------- SECTION ----------------------------
\section*{Step 2 — Determine Who or What Is Affected}
\paragraph{\textit{What are we looking for?} The individuals, groups, assets, and interests that could be harmed if the identified hazards materialise.}
\paragraph{\textit{Key techniques:}}
\begin{itemize}
  \item{Stakeholder mapping.}
  \item{Asset inventory.}
  \item{Process and value chain mapping.}
  \item{Exposure pathway analysis.}
\end{itemize}

\paragraph{\textit{Key questions:}}
\begin{itemize}
  \item{Who could be harmed — employees, customers, communities, third parties?}
  \item{What assets could be damaged — physical, financial, informational, reputational?}
  \item{How would harm be experienced — directly, indirectly, cumulatively?}
  \item{Are there particularly vulnerable parties who need specific consideration?}
\end{itemize}

%---------------------------- SECTION ----------------------------
\section*{Step 3 — Evaluate and Analyse Risks}
\paragraph{\textit{What are we looking for?} An assessment of the likelihood that the hazard will cause harm, and the severity of that harm if it occurs — taking into account existing controls.}
\paragraph{\textit{Key techniques:}}
\begin{itemize}
  \item{Risk matrix (likelihood × consequence).}
  \item{Bow-tie analysis.}
  \item{Fault tree analysis and event tree analysis.}
  \item{Quantitative risk assessment (where data supports it).}
  \item{Scenario analysis and stress testing.}
  \item{Expert judgement and structured elicitation.}
\end{itemize}

\paragraph{\textit{Key questions:}}
\begin{itemize}
  \item{How likely is the risk to materialise?}
  \item{How severe would the consequences be — across all relevant dimensions?}
  \item{How effective are existing controls?}
  \item{What is the residual risk after controls?}
  \item{Is the residual risk within appetite?}
\end{itemize}

\paragraph{\textit{Consequence dimensions to consider:}}
\begin{itemize}
  \item{Financial.}
  \item{Regulatory and legal.}
  \item{Operational.}
  \item{Reputational.}
  \item{Human (harm to people).}
  \item{Environmental.}
\end{itemize}

%---------------------------- SECTION ----------------------------
\section*{Step 4 — Decide on Controls and Actions}
\paragraph{\textit{What are we looking for?} The most appropriate and proportionate measures to reduce risk to an acceptable level — applying the Hierarchy of Controls.}
\paragraph{\textit{The Hierarchy of Controls (highest to lowest reliability):}}
\begin{itemize}
  \item{Elimination — remove the hazard entirely.}
  \item{Substitution — replace with a less hazardous alternative.}
  \item{Engineering controls — physical barriers and safeguards.}
  \item{Administrative controls — policies, procedures, and training.}
  \item{Personal protective equipment — last resort individual protection.}
\end{itemize}

\paragraph{\textit{Key questions:}}
\begin{itemize}
  \item{Can the risk be eliminated or substantially reduced at source?}
  \item{What controls would most effectively reduce likelihood?}
  \item{What controls would most effectively limit consequences?}
  \item{What detective controls are needed to identify the risk if it materialises?}
  \item{Are the proposed controls proportionate to the risk?}
  \item{Who is responsible for implementing each control.}
\end{itemize}


%---------------------------- SECTION ----------------------------
\section*{Step 5 — Record, Review, and Monitor}
\paragraph{\textit{What are we looking for?} Adequate documentation of the risk assessment, ongoing monitoring of the risk profile, and regular review to ensure the assessment remains current and the controls remain effective.}
\paragraph{\textit{Key requirements:}}
\begin{itemize}
  \item{Documented risk assessment to evidential standard.}
  \item{Key Risk Indicators defined, with thresholds and monitoring arrangements.}
  \item{Review schedule established — periodic and trigger-based.}
  \item{Ownership clearly assigned.}
  \item{Actions tracked with owners and deadlines.}
  \item{Changes to the risk profile incorporated promptly.}
\end{itemize}

\paragraph{\textit{Key questions:}}
\begin{itemize}
  \item{Is the documentation adequate to evidence the assessment?}
  \item{Are KRIs providing effective early warning?}
  \item{Have any trigger events occurred that require reassessment?}
  \item{Are controls being tested and confirmed as effective?}
  \item{Has the risk profile changed since the last assessment?}
\end{itemize}
